% Môn: Vật lý
% Lớp: 12
% Chương: Chương I. Vật lí nhiệt
% Bài: L12C1 Bài 2. Nội năng. Định luật I của nhiệt động lực học
% Loại: Dạng mẫu
% File riêng, không trộn vào de.tex
% Định nghĩa lệnh Lý thuyết — dùng khi biên dịch lt.tex bằng TeX.
% App web đọc lt.tex trực tiếp (bỏ \input file này).
% Trong lt.tex, từ thư mục bài:
%   % Định nghĩa lệnh Lý thuyết — dùng khi biên dịch lt.tex bằng TeX.
% App web đọc lt.tex trực tiếp (bỏ \input file này).
% Trong lt.tex, từ thư mục bài:
%   % Định nghĩa lệnh Lý thuyết — dùng khi biên dịch lt.tex bằng TeX.
% App web đọc lt.tex trực tiếp (bỏ \input file này).
% Trong lt.tex, từ thư mục bài:
%   \input{../../../../_lenh/lythuyet.tex}
% từ gốc repo:
%   \input{ngan-hang/_lenh/lythuyet.tex}
\makeatletter
\providecommand{\indam}[1]{\textbf{#1}}
\providecommand{\chuy}[1]{\par\smallskip\noindent\textit{#1}\par}
\providecommand{\luuy}[1]{\par\smallskip\noindent\textbf{Lưu ý.} #1\par}
\providecommand{\ghichu}[1]{\par\smallskip\noindent\textbf{Ghi chú.} #1\par}
\providecommand{\vidu}[1]{\par\smallskip\noindent\textbf{Ví dụ.} #1\par}
\providecommand{\Blythuyet}{\subsection*{Lý thuyết}}
% Hai cột: kiến thức | hình
\providecommand{\haicotchay}[2]{%
  \par\medskip\noindent
  \begin{minipage}[t]{0.52\linewidth}#1\end{minipage}\hfill
  \begin{minipage}[t]{0.46\linewidth}#2\end{minipage}%
  \par\medskip
}
\providecommand{\kienthuccot}[2]{\haicotchay{#1}{#2}}
% Dự phòng nếu chưa nạp graphicx / caption (không ghi đè bản thật)
\providecommand{\resizebox}[3]{#3}
\providecommand{\captionof}[2]{\par\smallskip\noindent\textit{#2}\par}
\newenvironment{hoatdong}[1]{%
  \par\medskip\noindent\textbf{Hoạt động.} #1\par\smallskip
}{\par\medskip}
\newenvironment{emcobiet}{%
  \par\medskip\noindent\textbf{Em có biết}\par\smallskip
}{\par\medskip}
\newenvironment{emdahoc}{%
  \par\medskip\noindent\textbf{Em đã học}\par\smallskip
}{\par\medskip}
\newenvironment{emcothe}{%
  \par\medskip\noindent\textbf{Em có thể}\par\smallskip
}{\par\medskip}
\newenvironment{kienthuc}{%
  \par\medskip\noindent\textbf{Kiến thức cốt lõi}\par\smallskip
}{\par\medskip}
\newenvironment{traloi}{%
  \par\smallskip\noindent\textit{Trả lời / ghi chú của em}\par
  \vspace{1.2em}%
}{\par\medskip}
\makeatother

% từ gốc repo:
%   % Định nghĩa lệnh Lý thuyết — dùng khi biên dịch lt.tex bằng TeX.
% App web đọc lt.tex trực tiếp (bỏ \input file này).
% Trong lt.tex, từ thư mục bài:
%   \input{../../../../_lenh/lythuyet.tex}
% từ gốc repo:
%   \input{ngan-hang/_lenh/lythuyet.tex}
\makeatletter
\providecommand{\indam}[1]{\textbf{#1}}
\providecommand{\chuy}[1]{\par\smallskip\noindent\textit{#1}\par}
\providecommand{\luuy}[1]{\par\smallskip\noindent\textbf{Lưu ý.} #1\par}
\providecommand{\ghichu}[1]{\par\smallskip\noindent\textbf{Ghi chú.} #1\par}
\providecommand{\vidu}[1]{\par\smallskip\noindent\textbf{Ví dụ.} #1\par}
\providecommand{\Blythuyet}{\subsection*{Lý thuyết}}
% Hai cột: kiến thức | hình
\providecommand{\haicotchay}[2]{%
  \par\medskip\noindent
  \begin{minipage}[t]{0.52\linewidth}#1\end{minipage}\hfill
  \begin{minipage}[t]{0.46\linewidth}#2\end{minipage}%
  \par\medskip
}
\providecommand{\kienthuccot}[2]{\haicotchay{#1}{#2}}
% Dự phòng nếu chưa nạp graphicx / caption (không ghi đè bản thật)
\providecommand{\resizebox}[3]{#3}
\providecommand{\captionof}[2]{\par\smallskip\noindent\textit{#2}\par}
\newenvironment{hoatdong}[1]{%
  \par\medskip\noindent\textbf{Hoạt động.} #1\par\smallskip
}{\par\medskip}
\newenvironment{emcobiet}{%
  \par\medskip\noindent\textbf{Em có biết}\par\smallskip
}{\par\medskip}
\newenvironment{emdahoc}{%
  \par\medskip\noindent\textbf{Em đã học}\par\smallskip
}{\par\medskip}
\newenvironment{emcothe}{%
  \par\medskip\noindent\textbf{Em có thể}\par\smallskip
}{\par\medskip}
\newenvironment{kienthuc}{%
  \par\medskip\noindent\textbf{Kiến thức cốt lõi}\par\smallskip
}{\par\medskip}
\newenvironment{traloi}{%
  \par\smallskip\noindent\textit{Trả lời / ghi chú của em}\par
  \vspace{1.2em}%
}{\par\medskip}
\makeatother

\makeatletter
\providecommand{\indam}[1]{\textbf{#1}}
\providecommand{\chuy}[1]{\par\smallskip\noindent\textit{#1}\par}
\providecommand{\luuy}[1]{\par\smallskip\noindent\textbf{Lưu ý.} #1\par}
\providecommand{\ghichu}[1]{\par\smallskip\noindent\textbf{Ghi chú.} #1\par}
\providecommand{\vidu}[1]{\par\smallskip\noindent\textbf{Ví dụ.} #1\par}
\providecommand{\Blythuyet}{\subsection*{Lý thuyết}}
% Hai cột: kiến thức | hình
\providecommand{\haicotchay}[2]{%
  \par\medskip\noindent
  \begin{minipage}[t]{0.52\linewidth}#1\end{minipage}\hfill
  \begin{minipage}[t]{0.46\linewidth}#2\end{minipage}%
  \par\medskip
}
\providecommand{\kienthuccot}[2]{\haicotchay{#1}{#2}}
% Dự phòng nếu chưa nạp graphicx / caption (không ghi đè bản thật)
\providecommand{\resizebox}[3]{#3}
\providecommand{\captionof}[2]{\par\smallskip\noindent\textit{#2}\par}
\newenvironment{hoatdong}[1]{%
  \par\medskip\noindent\textbf{Hoạt động.} #1\par\smallskip
}{\par\medskip}
\newenvironment{emcobiet}{%
  \par\medskip\noindent\textbf{Em có biết}\par\smallskip
}{\par\medskip}
\newenvironment{emdahoc}{%
  \par\medskip\noindent\textbf{Em đã học}\par\smallskip
}{\par\medskip}
\newenvironment{emcothe}{%
  \par\medskip\noindent\textbf{Em có thể}\par\smallskip
}{\par\medskip}
\newenvironment{kienthuc}{%
  \par\medskip\noindent\textbf{Kiến thức cốt lõi}\par\smallskip
}{\par\medskip}
\newenvironment{traloi}{%
  \par\smallskip\noindent\textit{Trả lời / ghi chú của em}\par
  \vspace{1.2em}%
}{\par\medskip}
\makeatother

% từ gốc repo:
%   % Định nghĩa lệnh Lý thuyết — dùng khi biên dịch lt.tex bằng TeX.
% App web đọc lt.tex trực tiếp (bỏ \input file này).
% Trong lt.tex, từ thư mục bài:
%   % Định nghĩa lệnh Lý thuyết — dùng khi biên dịch lt.tex bằng TeX.
% App web đọc lt.tex trực tiếp (bỏ \input file này).
% Trong lt.tex, từ thư mục bài:
%   \input{../../../../_lenh/lythuyet.tex}
% từ gốc repo:
%   \input{ngan-hang/_lenh/lythuyet.tex}
\makeatletter
\providecommand{\indam}[1]{\textbf{#1}}
\providecommand{\chuy}[1]{\par\smallskip\noindent\textit{#1}\par}
\providecommand{\luuy}[1]{\par\smallskip\noindent\textbf{Lưu ý.} #1\par}
\providecommand{\ghichu}[1]{\par\smallskip\noindent\textbf{Ghi chú.} #1\par}
\providecommand{\vidu}[1]{\par\smallskip\noindent\textbf{Ví dụ.} #1\par}
\providecommand{\Blythuyet}{\subsection*{Lý thuyết}}
% Hai cột: kiến thức | hình
\providecommand{\haicotchay}[2]{%
  \par\medskip\noindent
  \begin{minipage}[t]{0.52\linewidth}#1\end{minipage}\hfill
  \begin{minipage}[t]{0.46\linewidth}#2\end{minipage}%
  \par\medskip
}
\providecommand{\kienthuccot}[2]{\haicotchay{#1}{#2}}
% Dự phòng nếu chưa nạp graphicx / caption (không ghi đè bản thật)
\providecommand{\resizebox}[3]{#3}
\providecommand{\captionof}[2]{\par\smallskip\noindent\textit{#2}\par}
\newenvironment{hoatdong}[1]{%
  \par\medskip\noindent\textbf{Hoạt động.} #1\par\smallskip
}{\par\medskip}
\newenvironment{emcobiet}{%
  \par\medskip\noindent\textbf{Em có biết}\par\smallskip
}{\par\medskip}
\newenvironment{emdahoc}{%
  \par\medskip\noindent\textbf{Em đã học}\par\smallskip
}{\par\medskip}
\newenvironment{emcothe}{%
  \par\medskip\noindent\textbf{Em có thể}\par\smallskip
}{\par\medskip}
\newenvironment{kienthuc}{%
  \par\medskip\noindent\textbf{Kiến thức cốt lõi}\par\smallskip
}{\par\medskip}
\newenvironment{traloi}{%
  \par\smallskip\noindent\textit{Trả lời / ghi chú của em}\par
  \vspace{1.2em}%
}{\par\medskip}
\makeatother

% từ gốc repo:
%   % Định nghĩa lệnh Lý thuyết — dùng khi biên dịch lt.tex bằng TeX.
% App web đọc lt.tex trực tiếp (bỏ \input file này).
% Trong lt.tex, từ thư mục bài:
%   \input{../../../../_lenh/lythuyet.tex}
% từ gốc repo:
%   \input{ngan-hang/_lenh/lythuyet.tex}
\makeatletter
\providecommand{\indam}[1]{\textbf{#1}}
\providecommand{\chuy}[1]{\par\smallskip\noindent\textit{#1}\par}
\providecommand{\luuy}[1]{\par\smallskip\noindent\textbf{Lưu ý.} #1\par}
\providecommand{\ghichu}[1]{\par\smallskip\noindent\textbf{Ghi chú.} #1\par}
\providecommand{\vidu}[1]{\par\smallskip\noindent\textbf{Ví dụ.} #1\par}
\providecommand{\Blythuyet}{\subsection*{Lý thuyết}}
% Hai cột: kiến thức | hình
\providecommand{\haicotchay}[2]{%
  \par\medskip\noindent
  \begin{minipage}[t]{0.52\linewidth}#1\end{minipage}\hfill
  \begin{minipage}[t]{0.46\linewidth}#2\end{minipage}%
  \par\medskip
}
\providecommand{\kienthuccot}[2]{\haicotchay{#1}{#2}}
% Dự phòng nếu chưa nạp graphicx / caption (không ghi đè bản thật)
\providecommand{\resizebox}[3]{#3}
\providecommand{\captionof}[2]{\par\smallskip\noindent\textit{#2}\par}
\newenvironment{hoatdong}[1]{%
  \par\medskip\noindent\textbf{Hoạt động.} #1\par\smallskip
}{\par\medskip}
\newenvironment{emcobiet}{%
  \par\medskip\noindent\textbf{Em có biết}\par\smallskip
}{\par\medskip}
\newenvironment{emdahoc}{%
  \par\medskip\noindent\textbf{Em đã học}\par\smallskip
}{\par\medskip}
\newenvironment{emcothe}{%
  \par\medskip\noindent\textbf{Em có thể}\par\smallskip
}{\par\medskip}
\newenvironment{kienthuc}{%
  \par\medskip\noindent\textbf{Kiến thức cốt lõi}\par\smallskip
}{\par\medskip}
\newenvironment{traloi}{%
  \par\smallskip\noindent\textit{Trả lời / ghi chú của em}\par
  \vspace{1.2em}%
}{\par\medskip}
\makeatother

\makeatletter
\providecommand{\indam}[1]{\textbf{#1}}
\providecommand{\chuy}[1]{\par\smallskip\noindent\textit{#1}\par}
\providecommand{\luuy}[1]{\par\smallskip\noindent\textbf{Lưu ý.} #1\par}
\providecommand{\ghichu}[1]{\par\smallskip\noindent\textbf{Ghi chú.} #1\par}
\providecommand{\vidu}[1]{\par\smallskip\noindent\textbf{Ví dụ.} #1\par}
\providecommand{\Blythuyet}{\subsection*{Lý thuyết}}
% Hai cột: kiến thức | hình
\providecommand{\haicotchay}[2]{%
  \par\medskip\noindent
  \begin{minipage}[t]{0.52\linewidth}#1\end{minipage}\hfill
  \begin{minipage}[t]{0.46\linewidth}#2\end{minipage}%
  \par\medskip
}
\providecommand{\kienthuccot}[2]{\haicotchay{#1}{#2}}
% Dự phòng nếu chưa nạp graphicx / caption (không ghi đè bản thật)
\providecommand{\resizebox}[3]{#3}
\providecommand{\captionof}[2]{\par\smallskip\noindent\textit{#2}\par}
\newenvironment{hoatdong}[1]{%
  \par\medskip\noindent\textbf{Hoạt động.} #1\par\smallskip
}{\par\medskip}
\newenvironment{emcobiet}{%
  \par\medskip\noindent\textbf{Em có biết}\par\smallskip
}{\par\medskip}
\newenvironment{emdahoc}{%
  \par\medskip\noindent\textbf{Em đã học}\par\smallskip
}{\par\medskip}
\newenvironment{emcothe}{%
  \par\medskip\noindent\textbf{Em có thể}\par\smallskip
}{\par\medskip}
\newenvironment{kienthuc}{%
  \par\medskip\noindent\textbf{Kiến thức cốt lõi}\par\smallskip
}{\par\medskip}
\newenvironment{traloi}{%
  \par\smallskip\noindent\textit{Trả lời / ghi chú của em}\par
  \vspace{1.2em}%
}{\par\medskip}
\makeatother

\makeatletter
\providecommand{\indam}[1]{\textbf{#1}}
\providecommand{\chuy}[1]{\par\smallskip\noindent\textit{#1}\par}
\providecommand{\luuy}[1]{\par\smallskip\noindent\textbf{Lưu ý.} #1\par}
\providecommand{\ghichu}[1]{\par\smallskip\noindent\textbf{Ghi chú.} #1\par}
\providecommand{\vidu}[1]{\par\smallskip\noindent\textbf{Ví dụ.} #1\par}
\providecommand{\Blythuyet}{\subsection*{Lý thuyết}}
% Hai cột: kiến thức | hình
\providecommand{\haicotchay}[2]{%
  \par\medskip\noindent
  \begin{minipage}[t]{0.52\linewidth}#1\end{minipage}\hfill
  \begin{minipage}[t]{0.46\linewidth}#2\end{minipage}%
  \par\medskip
}
\providecommand{\kienthuccot}[2]{\haicotchay{#1}{#2}}
% Dự phòng nếu chưa nạp graphicx / caption (không ghi đè bản thật)
\providecommand{\resizebox}[3]{#3}
\providecommand{\captionof}[2]{\par\smallskip\noindent\textit{#2}\par}
\newenvironment{hoatdong}[1]{%
  \par\medskip\noindent\textbf{Hoạt động.} #1\par\smallskip
}{\par\medskip}
\newenvironment{emcobiet}{%
  \par\medskip\noindent\textbf{Em có biết}\par\smallskip
}{\par\medskip}
\newenvironment{emdahoc}{%
  \par\medskip\noindent\textbf{Em đã học}\par\smallskip
}{\par\medskip}
\newenvironment{emcothe}{%
  \par\medskip\noindent\textbf{Em có thể}\par\smallskip
}{\par\medskip}
\newenvironment{kienthuc}{%
  \par\medskip\noindent\textbf{Kiến thức cốt lõi}\par\smallskip
}{\par\medskip}
\newenvironment{traloi}{%
  \par\smallskip\noindent\textit{Trả lời / ghi chú của em}\par
  \vspace{1.2em}%
}{\par\medskip}
\makeatother


\subsection{Dạng bài tập mẫu}

\subsubsection{Dạng: Khái niệm nội năng và các yếu tố ảnh hưởng đến nội năng}
\begin{phuongphap}
\begin{enumerate}
    \item \textbf{Đọc và hiểu khái niệm nội năng:} Nội năng của một vật là tổng động năng và thế năng của các phân tử cấu tạo nên vật đó.
    \item \textbf{Xác định các yếu tố ảnh hưởng đến nội năng:}
    \begin{itemize}
        \item \textbf{Nhiệt độ:} Khi nhiệt độ tăng, động năng trung bình của các phân tử tăng, do đó nội năng tăng.
        \item \textbf{Thể tích (trạng thái):} Đối với chất khí, khi thể tích thay đổi (ví dụ, nén hoặc giãn nở), thế năng tương tác giữa các phân tử có thể thay đổi, dẫn đến thay đổi nội năng. Đối với chất rắn, lỏng, sự thay đổi thể tích ít ảnh hưởng đến nội năng hơn so với nhiệt độ.
        \item \textbf{Khối lượng/Số mol:} Vật có khối lượng lớn hơn (nhiều phân tử hơn) sẽ có nội năng lớn hơn ở cùng điều kiện nhiệt độ và thể tích.
    \end{itemize}
    \item \textbf{Phân tích các phát biểu:} Dựa vào khái niệm và các yếu tố ảnh hưởng, đánh giá tính đúng sai của từng phát biểu trong đề bài.
\end{enumerate}
\end{phuongphuap}
\begin{vidumau}
Xét các phát biểu sau về dạng “Khái niệm nội năng và các yếu tố ảnh hưởng đến nội năng”
\choice
{\True Nội năng của một vật phụ thuộc vào nhiệt độ và thể tích của vật.}
{Nội năng của một vật chỉ phụ thuộc vào nhiệt độ của vật.}
{Nội năng của một vật là tổng động năng của các phân tử cấu tạo nên vật.}
{Nội năng của một vật không thể thay đổi nếu không có sự truyền nhiệt.}
\loigiai{Nội năng của một vật là tổng động năng và thế năng của các phân tử cấu tạo nên vật. Nội năng phụ thuộc vào nhiệt độ (ảnh hưởng đến động năng) và thể tích (ảnh hưởng đến thế năng tương tác giữa các phân tử, đặc biệt với khí). Do đó, phát biểu A là đúng.}
\end{vidumau}

\subsubsection{Dạng: Các cách làm biến đổi nội năng của một vật}
\begin{phuongphap}
\begin{enumerate}
    \item \textbf{Nhận diện hai cách cơ bản:} Có hai cách để làm biến đổi nội năng của một vật:
    \begin{itemize}
        \item \textbf{Thực hiện công:} Khi có lực tác dụng làm vật dịch chuyển hoặc biến dạng, công được thực hiện và có thể làm thay đổi nội năng của vật (ví dụ: ma sát, nén khí).
        \item \textbf{Truyền nhiệt:} Khi có sự chênh lệch nhiệt độ giữa vật và môi trường xung quanh, nhiệt lượng sẽ truyền từ nơi có nhiệt độ cao hơn sang nơi có nhiệt độ thấp hơn, làm thay đổi nội năng của vật.
    \end{itemize}
    \item \textbf{Phân biệt hai cách:}
    \begin{itemize}
        \item \textbf{Thực hiện công:} Liên quan đến sự chuyển động vĩ mô của vật hoặc các phần của vật dưới tác dụng của lực.
        \item \textbf{Truyền nhiệt:} Liên quan đến sự trao đổi năng lượng ở cấp độ vi mô (chuyển động hỗn loạn của phân tử) do chênh lệch nhiệt độ.
    \end{itemize}
    \item \textbf{Áp dụng vào tình huống cụ thể:} Đọc kỹ đề bài, xác định hành động hoặc quá trình được mô tả và phân loại nó thuộc cách thực hiện công hay truyền nhiệt.
\end{enumerate}
\end{phuongphap}
\begin{vidumau}
Xét các phát biểu sau về dạng “Các cách làm biến đổi nội năng của một vật”
\choice
{\True Cọ xát một miếng kim loại làm nó nóng lên là một cách làm biến đổi nội năng bằng thực hiện công.}
{Đun nóng nước trong ấm là cách làm biến đổi nội năng bằng thực hiện công.}
{Để một vật nóng trong không khí, vật nguội đi là cách làm biến đổi nội năng bằng thực hiện công.}
{Nội năng của vật chỉ có thể thay đổi bằng cách truyền nhiệt.}
\loigiai{Cọ xát là quá trình có lực ma sát tác dụng và làm vật dịch chuyển tương đối, sinh công và chuyển hóa thành nội năng, làm vật nóng lên. Đây là cách biến đổi nội năng bằng thực hiện công. Đun nóng nước và vật nguội đi là các quá trình truyền nhiệt.}
\end{vidumau}

\subsubsection{Dạng: Lý thuyết về nội năng và các cách biến đổi nội năng}
\begin{phuongphap}
\begin{enumerate}
    \item \textbf{Tổng hợp kiến thức:} Kết hợp kiến thức về khái niệm nội năng, các yếu tố ảnh hưởng và hai cách biến đổi nội năng (thực hiện công và truyền nhiệt).
    \item \textbf{Đánh giá từng phát biểu:} Đọc từng phát biểu trong đề bài và so sánh với các kiến thức đã tổng hợp.
    \begin{itemize}
        \item Nội năng là hàm trạng thái, chỉ phụ thuộc vào trạng thái của hệ (nhiệt độ, thể tích).
        \item Nội năng có thể tăng hoặc giảm.
        \item Hai cách biến đổi nội năng là thực hiện công và truyền nhiệt.
        \item Công và nhiệt lượng là các đại lượng đặc trưng cho quá trình, không phải hàm trạng thái.
    \end{itemize}
    \item \textbf{Lưu ý các bẫy:} Cẩn thận với các phát biểu tuyệt đối (chỉ, luôn luôn, không thể) hoặc nhầm lẫn giữa nội năng với công/nhiệt lượng.
\end{enumerate}
\end{phuongphap}
\begin{vidumau}
Phát biểu nào sau đây đúng?
\choice
{\True Nội năng của một vật có thể thay đổi bằng cách thực hiện công hoặc truyền nhiệt.}
{Nội năng của một vật chỉ phụ thuộc vào nhiệt độ của vật.}
{Nhiệt lượng là một dạng năng lượng tồn tại trong vật.}
{Công là một dạng năng lượng tồn tại trong vật.}
\loigiai{Nội năng của một vật có thể thay đổi bằng hai cách: thực hiện công hoặc truyền nhiệt. Phát biểu B sai vì nội năng còn phụ thuộc vào thể tích. Phát biểu C và D sai vì nhiệt lượng và công là các quá trình trao đổi năng lượng, không phải dạng năng lượng tồn tại trong vật.}
\end{vidumau}

\subsubsection{Dạng: Nội năng và truyền nhiệt}
\begin{phuongphap}
\begin{enumerate}
    \item \textbf{Nhắc lại khái niệm truyền nhiệt:} Truyền nhiệt là quá trình trao đổi năng lượng giữa các vật hoặc giữa vật với môi trường do sự chênh lệch nhiệt độ.
    \item \textbf{Xác định chiều truyền nhiệt:} Nhiệt lượng luôn tự truyền từ vật có nhiệt độ cao hơn sang vật có nhiệt độ thấp hơn cho đến khi đạt trạng thái cân bằng nhiệt.
    \item \textbf{Mối liên hệ với nội năng:}
    \begin{itemize}
        \item Vật nhận nhiệt lượng thì nội năng tăng.
        \item Vật tỏa nhiệt lượng thì nội năng giảm.
    \end{itemize}
    \item \textbf{Phân tích tình huống:} Đọc kỹ đề bài để xác định các vật, nhiệt độ ban đầu của chúng và điều kiện tiếp xúc. Từ đó suy ra chiều truyền nhiệt và sự thay đổi nội năng.
\end{enumerate}
\end{phuongphap}
\begin{vidumau}
Cho hai vật $A$ và $B$ tiếp xúc nhau. Nhiệt chỉ tự truyền từ vật $A$ sang vật $B$ khi
\choice
{\True nhiệt độ của vật $A$ cao hơn nhiệt độ của vật $B$.}
{nội năng của vật $A$ lớn hơn nội năng của vật $B$.}
{khối lượng của vật $A$ lớn hơn khối lượng của vật $B$.}
{thể tích của vật $A$ lớn hơn thể tích của vật $B$.}
\loigiai{Nhiệt lượng luôn tự truyền từ vật có nhiệt độ cao hơn sang vật có nhiệt độ thấp hơn. Các yếu tố như nội năng, khối lượng hay thể tích không quyết định chiều truyền nhiệt.}
\end{vidumau}

\subsubsection{Dạng: Lý thuyết định luật I và ứng dụng}
\begin{phuongphap}
\begin{enumerate}
    \item \textbf{Nắm vững nội dung định luật I của nhiệt động lực học:} Độ biến thiên nội năng của một hệ bằng tổng nhiệt lượng mà hệ nhận được và công mà hệ nhận được.
    \item \textbf{Công thức:} $\Delta U = Q + A$
    \item \textbf{Quy ước dấu:}
    \begin{itemize}
        \item $\Delta U > 0$: Nội năng tăng.
        \item $\Delta U < 0$: Nội năng giảm.
        \item $Q > 0$: Hệ nhận nhiệt lượng.
        \item $Q < 0$: Hệ tỏa nhiệt lượng.
        \item $A > 0$: Hệ nhận công (môi trường thực hiện công lên hệ).
        \item $A < 0$: Hệ thực hiện công (hệ sinh công lên môi trường).
    \end{itemize}
    \item \textbf{Ứng dụng:} Định luật I được dùng để giải thích các quá trình nhiệt động lực học, nguyên lý hoạt động của động cơ nhiệt, tủ lạnh, v.v.
    \item \textbf{Phân tích phát biểu:} Đọc kỹ từng phát biểu, áp dụng công thức và quy ước dấu để kiểm tra tính đúng sai.
\end{enumerate}
\end{phuongphap}
\begin{vidumau}
Xét các phát biểu sau về dạng “Vận dụng định luật $I$ giải thích hiện tượng và thiết bị”
\choice
{\True Trong quá trình đẳng tích, toàn bộ nhiệt lượng mà chất khí nhận được dùng để làm tăng nội năng của nó.}
{Trong quá trình đẳng nhiệt, toàn bộ nhiệt lượng mà chất khí nhận được dùng để làm tăng nội năng của nó.}
{Động cơ nhiệt hoạt động dựa trên nguyên tắc chuyển hóa hoàn toàn nhiệt năng thành cơ năng.}
{Định luật I nhiệt động lực học cho biết nhiệt lượng và công là hai dạng năng lượng có thể chuyển hóa cho nhau.}
\loigiai{Trong quá trình đẳng tích ($V = \text{const}$), công mà chất khí thực hiện là $A = 0$. Theo định luật I: $\Delta U = Q + A = Q + 0 = Q$. Vậy, toàn bộ nhiệt lượng mà chất khí nhận được dùng để làm tăng nội năng của nó. Phát biểu B sai vì trong quá trình đẳng nhiệt $\Delta U = 0$, nên $Q = -A$. Phát biểu C sai vì không thể chuyển hóa hoàn toàn nhiệt năng thành cơ năng. Phát biểu D sai vì nhiệt lượng và công là các quá trình trao đổi năng lượng, không phải dạng năng lượng.}
\end{vidumau}

\subsubsection{Dạng: Tính độ biến thiên nội năng trong quá trình nhiệt động lực học}
\begin{phuongphap}
\begin{enumerate}
    \item \textbf{Xác định các đại lượng đã biết:} Đọc kỹ đề bài để xác định nhiệt lượng $Q$ và công $A$ mà hệ trao đổi với môi trường.
    \item \textbf{Áp dụng quy ước dấu:}
    \begin{itemize}
        \item Nếu hệ nhận nhiệt lượng, $Q > 0$. Nếu hệ tỏa nhiệt lượng, $Q < 0$.
        \item Nếu hệ nhận công (môi trường thực hiện công lên hệ), $A > 0$. Nếu hệ thực hiện công (hệ sinh công lên môi trường), $A < 0$.
    \end{itemize}
    \item \textbf{Sử dụng công thức định luật I:} $\Delta U = Q + A$.
    \item \textbf{Thực hiện phép tính và kiểm tra đơn vị:} Đảm bảo tất cả các đại lượng đều ở cùng một đơn vị (thường là Joule - J).
    \item \textbf{Kiểm tra dấu của kết quả:} Dấu của $\Delta U$ cho biết nội năng tăng ($+$) hay giảm ($-$).
\end{enumerate}
\end{phuongphap}
\begin{vidumau}
Một hệ nhận công 200 J và tỏa nhiệt 50 J. Độ biến thiên nội năng của hệ là bao nhiêu?
\choice
{\True $150 \text{ J}$.}
{$250 \text{ J}$.}
{$-150 \text{ J}$.}
{$-250 \text{ J}$.}
\loigiai{Hệ nhận công 200 J, nên $A = +200 \text{ J}$. Hệ tỏa nhiệt 50 J, nên $Q = -50 \text{ J}$.
Áp dụng định luật I nhiệt động lực học: $\Delta U = Q + A = -50 \text{ J} + 200 \text{ J} = 150 \text{ J}$.}
\end{vidumau}

\subsubsection{Dạng: Chưa phân dạng}
\begin{phuongphap}
\begin{enumerate}
\item \textbf{Đọc kỹ đề bài:} Xác định rõ yêu cầu của bài toán và các thông tin đã cho.
\item \textbf{Xác định kiến thức liên quan:} Dựa vào nội dung bài học (nội năng, định luật I nhiệt động lực học, các cách biến đổi nội năng), tìm ra các khái niệm, công thức, hoặc nguyên lý phù hợp để giải quyết bài toán.
\item \textbf{Áp dụng công thức/nguyên lý:}
    \begin{itemize}
        \item Nếu là bài toán tính toán: Ghi rõ công thức, thay số liệu (chú ý đơn vị và quy ước dấu cho $Q$ và $A$).
        \item Nếu là bài toán lý thuyết: Phân tích từng lựa chọn hoặc phát biểu dựa trên kiến thức đã học.
    \end{itemize}
\item \textbf{Kiểm tra kết quả:} Đối chiếu kết quả với các lựa chọn (nếu là trắc nghiệm) hoặc kiểm tra tính hợp lý của lời giải (nếu là tự luận).
\end{enumerate}
\end{phuongphap}
\begin{vidumau}
Một chất lỏng có khối lượng $m$, nhiệt hoá hơi riêng là $L$, nhiệt nóng chảy riêng là $\lambda$, nhiệt dung riêng $c$. Nhiệt hóa hơi có thể được xác định thông qua công thức nào sau đây?
\choice
{\True $Q = mL$}
{$Q = mc\Delta T$}
{$Q = m\lambda$}
{$Q = L/m$}
\loigiai{Nhiệt hóa hơi là nhiệt lượng cần cung cấp để hóa hơi hoàn toàn một khối lượng chất lỏng ở nhiệt độ sôi. Công thức tính nhiệt lượng hóa hơi là $Q = mL$, trong đó $m$ là khối lượng chất lỏng và $L$ là nhiệt hóa hơi riêng.}
\end{vidumau}
